% ---------------------------------------------------------------------------
% Paste into chapters/04-design.tex (the data model section) and
% chapters/05-implementation.tex or 05-evaluation.tex (the results section).
% Needs \usepackage{booktabs} in main.tex. No other new packages.
% ---------------------------------------------------------------------------

\section{The data model}
\label{sec:datamodel}

A school timetable, as schools actually keep it, is a grid of two layers: a subject label and a
teacher label in each cell. The central finding of this phase of the project is that
\emph{the same visual pattern in that grid carries four different meanings}, and that nothing in
the file distinguishes them. A model that reads the grid under a single interpretation cannot
reproduce the school's own timetable. Before the four were separated, the real 2026--27 timetable
of KMBVM Pottore was \emph{infeasible} under our constraints in every band. After they were
separated, all four bands reproduce.

\subsection{Four patterns that look alike}

\paragraph{An in-section split.} One cell, several teacher codes, one section: the class divides
and goes to several rooms at once. \texttt{SCI LAB} under \texttt{UPR/REM/TSD} is physics,
chemistry and biology running simultaneously for thirds of one class. Every named teacher is
occupied for every lesson of that pair.

\paragraph{A regroup across sections.} The same teacher code appears in two sections in the same
period. Since a teacher cannot be in two rooms, those two groups are one class: the sections'
students regroup. This covers \texttt{DMD} (dance, music and drawing for two sections at a time
under three specialists), \texttt{M/IP/P} across 11E--11G, the languages of 10A, 10B and 10C
taken together, and 11A's biology students sitting in 11E's biology lesson. The teacher is
occupied once, not once per section. Groups are detected from the grid itself: a shared teacher
together with an overlapping subject name, treating a laboratory and its subject as the same
subject, so that \texttt{BIO LAB} matches \texttt{BIO}. That equivalence is what identifies the
cross-label cases.

\paragraph{A shared load.} One teacher per cell, but two teachers across the week. 12B English is
ANR five times and BIN once. This is not a split: one teacher takes each lesson and each
teacher's weekly share is fixed. It is modelled with a per-teacher choice variable beneath each
lesson variable, constrained to sum to the lesson and to meet each teacher's observed share.

\paragraph{A double-period continuation.} In grades 11--12 every laboratory is a double period,
written as a label followed by a blank cell. The blank is not free time: the class and its
teachers are still present. This must be filled in \emph{before} anything else is measured,
because regroups, teacher loads and conflicts are all counted from the expanded grid.

\subsection{What a shared teacher does not imply}

Two sections taught by the same teacher are not automatically one block. ANR teaches 12B and 12E
English together five times a week and separately once. They share a teacher; they are not tied
to a common period. Sections are therefore aligned in time only where their own grid never
separates them, rather than wherever their lesson counts happen to agree.

\subsection{Cells the model cannot explain}

The same reasoning detects errors in the source. A single appearance by a teacher who takes that
subject nowhere else in the band is read as a cell written in the wrong row rather than as a
teaching assignment; the lesson is given to the teacher who takes the rest, and the cell is
reported. Four were found this way, listed in Table~\ref{tab:wrongrow}. One further cell is
reported and remains unexplained: ANP appears in 6C Hindi and 8C English in the same period,
across different grades and unrelated subjects.

\begin{table}[ht]
\centering
\caption{Cells in which a teacher appears in a subject she teaches nowhere else.}
\label{tab:wrongrow}
\begin{tabular}{lll}
\toprule
Teacher & Appears in & Actually teaches \\
\midrule
ANR & 11B physics  & English \\
ADS & 10B English  & mathematics \\
ANP & 6C Hindi     & English \\
PMN & 6D English   & computer science / mathematics \\
\bottomrule
\end{tabular}
\end{table}

\subsection{Slots that do not exist}

Grades 1--5 open with a zero period on Thursday only, giving 46 slots in a week rather than 45.
The model carries a set of closed slots, declared for the whole school or for an individual
section, so that an irregular day requires no special case anywhere else in the formulation. In a
single whole-school solve the zero period is closed for every secondary section and open for the
primary ones.

% ---------------------------------------------------------------------------

\section{Evaluation}
\label{sec:evaluation}

\subsection{Reproduction as the test of a model}

The literature on automated timetabling generally reports that a feasible timetable was produced.
That demonstrates less than it appears to, because a cold solve conceals modelling error: a model
that misstates the school will still return \emph{some} timetable, and there is nothing to compare
it against. We therefore adopt a stronger and falsifiable test.

\begin{quote}
\emph{Is the school's own working timetable feasible under our constraints?}
\end{quote}

The school's timetable is known to work, since the school is running on it. If it is infeasible in
our model, the model asserts something false about the school, and the fault is ours. Every lesson
is pinned to the slot the school actually used and the model is asked whether that assignment
holds. Table~\ref{tab:reproduction} gives the result. Each of the four failures that this test
exposed corresponded to one of the four patterns of Section~\ref{sec:datamodel}; none of them
would have been visible from a cold solve alone.

\begin{table}[ht]
\centering
\caption{Reproduction: the school's own timetable, every lesson pinned where the school put it.}
\label{tab:reproduction}
\begin{tabular}{lrrl}
\toprule
Band & Sections & Lessons pinned & Result \\
\midrule
Grades 11--12 & 12 & 474 & feasible \\
Grades 9--10  &  8 & 360 & feasible \\
Grades 6--8   & 12 & 540 & feasible \\
Grades 1--5   & 18 & 828 & feasible \\
\bottomrule
\end{tabular}
\end{table}

\subsection{Cold start}

The second test is construction: given only the derived instance and the stated rules, with no
reference to the existing timetable, can the system build one? Bands are solved in the order the
school's own committee works in, each against the teacher time the earlier bands have already
consumed, since 30 teacher codes appear in both grades 6--8 and 9--10 and 11 appear in all three
secondary bands. Table~\ref{tab:coldstart} gives the result. An independent check of the output
found no section with the wrong lesson counts and no teacher in two places at once anywhere in the
school.

\begin{table}[ht]
\centering
\caption{Cold start: built from the instance alone, in the school's own build order.}
\label{tab:coldstart}
\begin{tabular}{llrr}
\toprule
Band & Result & Time (s) & Rules \\
\midrule
Grades 11--12 & feasible, 9 preferences relaxed & 110 & 51 \\
Grades 9--10  & optimal                         &  11 & 32 \\
Grades 6--8   & optimal                         & 108 & 24 \\
Grades 1--5   & feasible, 4 preferences relaxed & 110 & 26 \\
\bottomrule
\end{tabular}
\end{table}

Preferences that cannot be met are reported rather than silently discarded, which is the
behaviour the project argues for: the system states what it gave up.

\subsection{One pass against band by band}

Solving all 50 sections in a single pass returns \textsc{unknown} after five and a half minutes.
Fixing the band-by-band solution into that same one-pass model and asking whether it holds returns
\emph{optimal}, which establishes that the one-pass formulation is correct and merely hard to
search. Separating a modelling fault from a search cost in this way is only possible because the
reproduction test exists.

Obtaining that result also exposed a genuine fault in our own work, and the system rather than the
author found it. An early attempt placed the four bands on a common grid by starting the shorter
days later, which moved the lunch break relative to each band's own period numbering, so that a
laboratory double legal within its band straddled the break. Of the 100 rules in the model, the
conflict explainer named the single offending rule, \texttt{BIO/CS LAB is a double period}, in
0.2 seconds. Each band now retains its own period numbering, with surplus slots at the end of the
day closed instead.

\subsection{Discussion}

The measured result of this phase is that the dominant source of error was neither the solver nor
the natural-language layer. Solving was never the bottleneck, at 11 to 110 seconds per band, and
the language layer was in place from September. What made the system wrong about a real school was
the data model, and specifically four ambiguities that no one had written down, because the person
who builds the timetable has never needed to write them down. This locates the project's
contribution more precisely than the original framing did: the system must detect that it is
guessing not only about what the coordinator \emph{says}, but about what her existing timetable
\emph{already encodes}.
